\subsubsection{Independent paired worlds}

MatrixGroupWorld provides two independently resampled worlds, each with new data and a new initialization. Within each world, four task specialists share a saved, genuine initialization and each completes the fixed 20,000-update budget. We report both worlds separately, with equal weights over $n\in\{6,8,10\}$ and the same prespecified transformations, controls, and metrics. These worlds repeat the matrix setting; they do not independently confirm PermWorld across initializations. Because data and initialization vary together, differences between worlds do not isolate a causal effect of initialization. The four relations, three sizes, and three metrics are coupled measurements rather than additional independent replications.

\subsubsection{R1: an illustrative case}

R1 relates $F_1$ and $F_2$ through transposition. Both tasks achieve 100\% full-vocabulary answer accuracy in both worlds at every size on the complete identity training, validation, and analysis sets. The $F_2$ models also achieve 100\% analysis accuracy under the correct and wrong R1 transformations. Nevertheless, their geometric pairing preference depends on the control and metric (Figure~\ref{fig:matrixcase}; Table~\ref{tab:matrix_r1}). This makes R1 an illustrative case, rather than evidence for a general reversal of relation preference.

For raw, size-stratified hidden-state CKA, worlds one and two respectively give $(B_t,\Delta)=(0.520209,0.519581)$ and $(0.558791,0.553006)$. After joint-answer-stratum control, these become $(-0.019384,-0.020747)$ and $(-0.014022,-0.014551)$. Thus, these controlled, equally weighted CKA estimates favor the wrong pairing. Answer control changes both centering and retained sample support; this contrast cannot be interpreted as a causal subtraction of all answer information.

The fixed-transform comparison supplies an additional distinction. For R1,
the designated task ranks uniquely first among three candidates under answer
control in both worlds, both initially and after training. Its advantage over
the two non-designated tasks changes from $0.1894$ to $0.1182$ and from
$0.1179$ to $0.0663$. Thus first rank in Comparison B can coexist with a
negative input-pairing preference in Comparison A, and its controlled task
advantage need not originate at training.

The controlled sensitivity metrics give different aggregate directions. For $k$NN overlap, $(B_t,\Delta)$ is $(0.047718,0.039648)$ and $(0.081837,0.069883)$; for held-out Procrustes it is $(0.052454,0.038940)$ and $(0.045210,0.036150)$. Procrustes is scored as negative alignment error, so larger values indicate better alignment. These metrics use different scales; their numerical magnitudes are not directly comparable. Raw $B_t$ and $\Delta$ are positive for all three metrics in both worlds, whereas controlled CKA disagrees with the two sensitivity metrics. This aggregate disagreement also has size-specific qualifications: controlled CKA has positive $B_t$ and $\Delta$ at $n=10$ in both worlds, while controlled Procrustes has negative $B_t$ and $\Delta$ at $n=6$ in both. World one's controlled $k$NN $B_t$ and $\Delta$ are negative at $n=8$. We retain all sizes rather than select a uniformly positive or negative subset.

R3 illustrates why $\Delta$ must also be decomposed. Its controlled CKA increments are $0.048498$ and $0.121332$, with correct-pair changes $\delta_c=-0.800773$ and $-0.755736$, and wrong-pair changes $\delta_u=-0.849270$ and $-0.877068$. Both scores decrease; the larger wrong-pair decrease produces the positive increment. These results describe a change in relative selectivity, not increased correct-pair convergence. Interpretation is further limited by weak $F_3$ learning at $n=10$: training accuracies are 45.61\% and 39.21\%, validation accuracies 40.23\% and 31.25\%, and identity-analysis accuracies 35.55\% and 31.64\%. Accuracy on the correct R3-transformed analysis inputs remains 35.74\% and 32.23\%. The low training accuracy prevents attributing this weakness solely to generalization.

We interpret these observations as measured pairing selectivity, with algebraic and causal mechanisms unresolved. Numerical checking retains its original verification caveat.\footnote{The unchanged audit \texttt{c2e799\ldots} records 64 passed, two failed, and six consistently undefined checks at tolerance $10^{-8}$. Both failures concern trained raw output-probability $k$NN, for correct and wrong pairings, on the fixed 64-input first-world R1/$n=6$ subset. They remain failed; see Appendix~\ref{app:verification}.}
